\scanpage{66}

that \deleted{$\unsure{M_0}$}
\inserted{the subgroup under consideration of $\Pi_0\subset\mathrm{Sl}(2,\mathbb Z)$,
of index $3$ in $\Pi_0$,}
is the inverse image of this preceding subgroup \deleted{in $M$} of
$\mathrm{Sl}(2,\mathbb Z/12\mathbb Z)$.

Let us sum up the results obtained:

\uline{Proposition} Let us consider in $\mathcal G_0^+=\mathrm{Sl}(2,\mathbb Z)$
the tower of invariant subgroups (even in $\mathcal G_0$)
\[
(114)\qquad
\diagram{$\begin{tikzcd}[column sep=small,row sep=small]
\mathcal G_0^+&
\Pi=\Pi[2]\arrow[l,hook',"6"']\arrow[d,equal,"\mathrm{dfn}"]&
\Pi_0\arrow[l,hook',"2"']\arrow[d,equal]&
\Pi_{00}\ (\supset[\Pi_0,\Pi_0])\arrow[l,hook',"3"']\arrow[d,equal]\\
&\begin{array}{c}
\{u\in\mathcal G_0^+\mid u\equiv1\ (2)\}\\
\deleted{\Vert}\\
\deleted{\unsure{\mathfrak{sl}(2,\mathbb F_2)}\simeq
\{u\in\mathrm M_2(\mathbb F_2)\mid\mathrm{Tr}\,u=0\}}\\
\unsure{(i)|(1)}
\end{array}
&\begin{array}{l}
\text{Subgroup}\\
\text{generated}\\
\text{by the }-\ell_i=\lambda_i
\end{array}
&\begin{array}{l}
\text{Subgroup }\uline{\text{invariant}}\\
\text{generated by}\\
\deleted{\text{by the }\lambda_i,}\\
\deleted{[\lambda_0,\lambda_1]\ \text{and}}\\
\lambda_1\lambda_0^{-1}=\ell_1\ell_0^{-1}
\end{array}
\end{tikzcd}$}
\]

Let us consider on the other hand the quotient group
\[
(115)\qquad
\diagram{\begin{tikzcd}[row sep=small,column sep=small]
\mathrm{Sl}(2,\mathbb Z/12\mathbb Z)
\arrow[r,phantom,"\simeq" description]
\arrow[d,phantom,"\simeq" description]&
\mathcal G_0^+/\underbrace{\Pi[12]}_{\simeq\{u\in\mathcal G_0^+\mid u\equiv1\ (12)\}}\\
\mathrm{Sl}(2,\mathbb Z/4\mathbb Z)\times
\mathrm{Sl}(\unsure{2},\mathbb Z/\unsure{3}\mathbb Z),&
\end{tikzcd}}
\]
and the tower of the following subgroups
\[
(116)\qquad
\diagram{$\left\{\begin{array}{rll}
&H=\mathrm{Sl}(2,\mathbb Z/4\mathbb Z)\times\mathrm{Sl}(2,\mathbb Z/3\mathbb Z)
&\text{order }48\times24\\
6\ \cup&\hfill\cup\hfill&\\
&\Pi^\bullet=\mathfrak{sl}(2,\mathbb F_2)\times\mathrm{Sl}(2,\mathbb Z/3\mathbb Z)
&\text{order }8\times24\\
2\ \cup&\hfill\cup\hfill&\\
&\Pi_0^\bullet=\mathfrak{sl}(2,\mathbb F_2)'\times\mathrm{Sl}(2,\mathbb Z/3\mathbb Z)
&\text{order }4\times24\\
3\ \cup&\hfill\cup\hfill&\\
&\Pi_{00}^\bullet=\mathfrak{sl}(2,\mathbb F_2)'\times
\underbrace{\mathrm{Syl}_2}_{\simeq\deleted{\illegible{1}}D_4}
&\text{order }4\times8,
\end{array}\right.$}
\]

\margin{\uline{N.B.} $(\mathcal G_0^+)_{\unsure{\mathrm{ab}}}\simeq
\mathbb Z/12\mathbb Z$ is identified with a
\deleted{$\unsure{(115)|(116)}\ \unsure{\mathrm{Sl}(2)}$}
quotient of $H/\Pi_{00}^\bullet$ by a subgroup (116) of order $3$
--- that is again the subgroup generated by the element
$(-\dot\rho,1)$ of
\[
\underbrace{\mathrm{Sl}(2,\mathbb Z/4\mathbb Z)/\mathfrak{sl}(2,\mathbb F_2)'}_{\mathcal G_0^+/\Pi_0}
\times\mathrm{Sl}(2,\mathbb Z/3\mathbb Z)/\mathrm{Syl}_2.
\]}

Where \deleted{\unsure{we have}}
\[
(117)\qquad
\diagram{$\left\{\begin{array}{ll}
\text{a)}&\mathfrak{sl}(2,\mathbb F_2)\simeq
\mathrm{Ker}\bigl(\mathrm{Sl}(2,\mathbb Z/4\mathbb Z)\longrightarrow
\mathrm{Sl}(2,\mathbb Z/2\mathbb Z)\bigr)\\[.5em]
\text{b)}&\mathfrak{sl}(2,\mathbb F_2)'\unsure{\subset}\mathfrak{sl}(2,\mathbb F_2),
\quad\begin{array}{l}
\text{normal subgroup }\simeq\\
\deleted{\unsure{\mathbb F_3|\mathbb F_2}}\ \mathbb F_2\times\mathbb F_2\ ;\
\text{whose three elements }\ne0\\
\text{are }\dot\lambda_0,\dot\lambda_1,\dot\lambda_\infty\
\text{ (images}\\
\text{in }\mathrm{Sl}(2,\mathbb Z/4\mathbb Z)\text{) --- it is}\\
\text{a complement to its}\\
\text{Lie subalgebra }\mu\simeq\mathbb F_2\ \text{of matrices}\\
\text{\unsure{constant|scalar}}
\end{array}\\[.5em]
\text{c)}&\underset{\substack{\Vert\\
\mathrm{Ker}(\mathrm{Sl}(2,\mathbb Z/3\mathbb Z)\to\mathcal G_3^+)}}{\mathrm{Syl}_2}
\subset\mathrm{Sl}(2,\mathbb Z/3\mathbb Z),\quad
\begin{array}{l}
\text{Sylow subgroup, of order }8\\
\text{(which is normal)}
\end{array}
\end{array}\right.$}
\]

\scanpage{67}
\mspage{681}

This being so, the tower of subgroups (114) is the inverse image of the tower (116).

\deleted{\unsure{Remark|Remark}}
To satisfy my conscience, I am going to take up the preceding \unsure{sketch|sketch}
of a proof again. We note that the homomorphism
$\mathcal G_0^+\longrightarrow H$ maps the groups appearing in
(114) into those of (116) --- for $\Pi$ and $\Pi_0$, we have
in fact known for a while that they are the inverse images of
$\Pi^\bullet$, $\Pi_0^\bullet$! $\unsure{\ast|\leftarrow}$
(For this, it \unsure{was enough|would be enough} to examine
$\mathrm{Sl}(2,\mathbb Z/4\mathbb Z)$, instead of
$\mathrm{Sl}(2,\mathbb Z/12\mathbb Z)$ \ldots).

Thus the inverse image of $\Pi_{00}^\bullet$ in $\mathcal G_0^+$
contains $\Pi_{00}$. As both are of index $36$, these two
groups are equal, q.e.d. (The same
\unsure{argument} could be applied to the case of $\Pi_0$, $\Pi_0^\bullet$ \ldots)

\uline{Corollary} The canonical homomorphism
\[
\diagram{\begin{tikzcd}[column sep=small]
\mathcal G_0^*=\mathrm{Gl}(2,\mathbb Z)\arrow[r]&
\mathrm{Gl}(2,\mathbb Z/12\mathbb Z)/
\bigl(\Pi_{00}^\bullet\deleted{\unsure{\simeq}}=
\mathfrak{sl}(2,\mathbb F_2)'\times\mathrm{Syl}_2\bigr)
\end{tikzcd}}
\]
\unsure{factors} through
\[
\mathcal G=\deleted{\mathcal G_0}\mathrm{Gl}(2,\mathbb Z)/[\Pi_0,\Pi_0]
\]
studied in the present \S\,51, and the inverse image of
$\Pi_{00}^\bullet$ is none other than the subgroup $M_0$ of (99) (108).

\margin{\uline{N.B.}
But $\mathcal G_0^{\unsure{+|*}}\longrightarrow
\unsure{\mathrm{Sl}|\mathrm{Gl}}(2,\mathbb Z/12\mathbb Z)$
does \uline{not} factor through $\mathcal G^{\unsure{+|*}}$,
for if this were so $\Pi_0^{\unsure{\bullet|\prime}}$
would be a \unsure{sub-}quotient of $\Pi_0/[\Pi_0,\Pi_0]$ and
hence would be commutative --- \unsure{now all}
\unsure{we see} \unsure{two|therefore} \unsure{which} \unsure{are}
both.}

\scanpage{68}

The preceding considerations make fairly \unsure{natural|natural} the consideration of the covering of order $3$ of $U_{0,\overline{\mathbb Q}}$ \inserted{\deleted{\unsure{or of} $M_{1,\unsure{1|0}}$}} corresponding to the subgroup $\Pi_{00}$ of $\Pi_0\simeq\Pi_1(U_{0,\overline{\mathbb Q}})$.

\margin{it is even a \unsure{covering} of $M_{1,1}$\dots}

We check at once that it is a curve \uline{of genus $1$} -- if then we take it into our heads to work with $[\Pi_{0,0},\Pi_{0,0}]$ instead of $[\Pi_0,\Pi_0]$, we find a group extension of $\Pi_0/[\Pi_0,\Pi_0]\simeq M$, by a group of rank $2$ which is ``elliptic'', and on which $M$ will operate in a highly nontrivial way -- later on we would have to examine this action, and put it in relation with the action we have \unsure{found}, subordinate to the homomorphism
\diagram{\begin{tikzcd}
\mathrm{Sl}(2,\mathbb Z)^{\wedge}\arrow[r] & \mathrm{Sl}(2,\widehat{\mathbb Z}).
\end{tikzcd}}

\margin{it is an elliptic curve of invariant $0$ (for it is equipped with an \unsure{automorphism} of order $3$)}

\scanpage{69}
\mspage{682}

5) The groups $\underline M_{\rho,\sigma}$, $\underline M_{\rho,\sigma}^{\sharp}$ (via $\underline N_{\rho,\sigma}$).

To every $\ell\in\Gamma\unsure{\underline M_0}$ there corresponds an inner automorphism $\operatorname{int}(\ell)$, which is the identity on $\underline M_0$. We will have, tautologically
\[
(118)\qquad
\left\{\begin{aligned}
\operatorname{int}(\ell)(\rho)&=\alpha_{\unsure{0|\cdot}}(\rho)\\
\operatorname{int}(\ell)(\sigma)&=\beta(\sigma)
\end{aligned}\right.
\qquad\alpha=\beta=\ell
\]

But \deleted{we will have} will we for all that have the loop relation (relative to $\underline\mu=1$ hence $\nu''=0$)
\[
(1-\rho_M)(\underbrace{\alpha}_{\ell})=(1-\sigma_M)(\underbrace{\beta}_{\ell}).
\]
\[
\underset{\unsure{(*)|\alpha_1}}{\text{i.e.}}\qquad
\rho_M(\ell)=\sigma_M(\ell)\qquad ?
\]

Setting again
\[
(119)\qquad\ell=c_0\lambda_0+c_1\lambda_1
\]
\deleted{we will have}
the relation $(*)$ is written
\[
c_0\lambda_1+c_1\underbrace{\lambda_\infty}_{-\lambda_0-\lambda_1}
=c_0\lambda_1+c_1\lambda_0
\]
\[
\text{i.e.}\qquad-\unsure{c_1}\lambda_0+(c_0-c_1)\lambda_1
=c_1\lambda_0+c_0\lambda_1
\]
i.e. $c_1=0$ i.e. $\ell=c_0\lambda_0$. Thus it is subject to the \unsure{condition|conditions} that we restrict ourselves to considering $\ell\in\Gamma\underline L_0^{\unsure{\underline M_0|\underline M}}\unsure{\Rightarrow}$; \unsure{corresponding|corresponding}

\diagram{\begin{tikzcd}[column sep=small]
{(120)\quad\deleted{\unsure{\underline L_0}}\quad\ell=c_0\lambda_0\in
\underline L_{0\deleted{\unsure{,\rho,\sigma}}}^{\underline M}
\overset{\mathrm{dfn}}{=}\operatorname{Im}(}
& \underline{\mathbb Z}\arrow[r]\arrow[d,phantom,"\simeq" description]
& \underline M)\\
& \mathbb G_a\arrow[d,phantom,"\unsure{\ni}" description] &\\
& c_0\arrow[r,mapsto] & c_0\lambda_0
\end{tikzcd}}

We have a \deleted{\unsure{dia}} diagram
\diagram{\begin{tikzcd}[column sep=large]
{(121)\quad\mathbb G_a\simeq\underline L_0^{\unsure{\underline M_0|\underline M}}}
\arrow[r]\arrow[d]
& \underline N_{\rho,\sigma} & \deleted{\unsure{\underline L_0^{\underline M}}}\\
\underline{\mathcal G}^{+} & &
\end{tikzcd}}

\scanpage{70}

the image of
\diagram{\begin{tikzcd}
{(\unsure{127|122})\quad\underline L_0^{\unsure{\underline M_0|\underline M}}}\arrow[r]
& \underline N_{\rho,\sigma}
\end{tikzcd}}
is invariant -- the homomorphism $(121)$ being in fact compatible with the action of $N_{\rho,\sigma}$ -- via $\underline\mu$ on $\underline L_0^{M_0}\simeq\mathbb G_a$ (i.e. induced by the operation of $\underline N_{\rho,\sigma}$ \unsure{of} on\unsure{!| the} $\underline{\mathcal G}^{+}\supset\underline M\supset\underline L_0^{\underline M}$), and \deleted{the}\inserted{the two} operations of \deleted{$\unsure{N_\rho|A_\rho}$} $\underline L_0^{\underline M}$ on $\underline{\mathcal G}^{+}$ deduced from $(121)$ being \unsure{identical|identical}. Thus according to the general \unsure{theory|theorem} we \unsure{find|will find} a group scheme $\underline M_{\rho,\sigma}$ and a diagram of exact sequences

\diagram{\begin{tikzcd}[column sep=large,row sep=large]
{(123)\quad1}\arrow[r]
&\underline L_0^{\underline M}\arrow[r]\arrow[d]
&\underline N_{\rho,\sigma}\arrow[r,Rightarrow]\arrow[d]
&\overset{\inserted{\overset{\mathrm{dfn}}{=}\underline N_{\rho,\sigma}/\underline L_0^{\underline M}}}{\underline\Gamma_{\rho,\sigma}}\arrow[r]\arrow[d,Rightarrow]
&1\\
1\arrow[r]
&\underline{\mathcal G}^{+\unsure{a|\alpha}}\arrow[r]
&\underline M_{\rho,\sigma}\arrow[r,Rightarrow]
&\underline\Gamma_{\rho,\sigma}\arrow[r]
&1
\end{tikzcd}}

\margin{with the notations of \S\,49 $\overline{\underline{\mathrm{VI}}}$, the group denoted here $\underline M_{\rho,\sigma}$ would probably be denoted $\underline M_{\rho,\sigma}/\unsure{H_\rho|H_{\underline\Pi}}$, and $\underline M_{\rho,\sigma}^{\sharp}$ would be $\underline M_{\rho,\sigma}^{\sharp}/\unsure{H_\rho}$ (loc. cit. $8^\circ$)}

Replacing $\underline{\mathcal G}^{+}$ by $\underline M$, we likewise find $\underline M_{\rho,\sigma}^{\sharp}$, and

\diagram{\begin{tikzcd}[column sep=large,row sep=large]
{(124)}\arrow[d,phantom,"\Bigg\{" description]
&1\arrow[r,dash]
&\underline L_0^{\underline M}\arrow[r,dash]\arrow[d]
&\underline N_{\rho,\sigma}\arrow[r,dash]\arrow[d]
&\underline\Gamma_{\rho,\sigma}\arrow[r,dash]\arrow[d,Rightarrow]
&1\\
{} &1\arrow[r]
&\deleted{\unsure{\underline{\mathcal G}^{+}}}\underline M\arrow[r]
&\underline M_{\rho,\sigma}^{\sharp}\arrow[r]
&\underline\Gamma_{\rho,\sigma}\arrow[r]
&1
\end{tikzcd}}
,

moreover we have a homomorphism from $(124)$ into $(123)$ (compare \S\,\uline{49} $\overline{\underline{\mathrm{VI}}}$ $5^\circ$\dots), and
\[
(125)\qquad M_{\rho,\sigma}^{\sharp}\subset M_{\rho,\sigma}\qquad\dots
\]

\scanpage{71}
\mspage{683}

We must specify the structure of $\Gamma_{\rho,\sigma}$, noting that $L_0^{\underline M}\subset N_{\rho,\sigma}^{0+}$, hence
\diagram{\begin{tikzcd}
{(\unsure{126|125})\quad1}\arrow[r]
&\deleted{\unsure{N_{\rho,\sigma}}}\underline\Gamma_{\rho,\sigma}^{0+}\arrow[r]
&\underline\Gamma_{\rho,\sigma}\arrow[r]
&\underline{\mathbb O}_{12}^{*}\arrow[r]
&1
\end{tikzcd}}
with
\diagram{\begin{tikzcd}
{(127)\quad1}\arrow[r]
&\underline L_0^{M}\arrow[r,thick]
&\underline N_{\rho,\sigma}^{0+}\arrow[r]
&\underline\Gamma_{\rho,\sigma}^{0+}\arrow[r]
&1\\
&&\deleted{\unsure{\text{it is }C_f|\text{it is }G_f}}&&
\end{tikzcd}}

Let us consider \deleted{the}\inserted{the}\deleted{exact sequence}\inserted{surjective homomorphism} of $\mathbb Z$-modules

\diagram{\begin{tikzcd}[column sep=small,row sep=small]
{(128)\quad\unsure{N_{0,\sigma}^{0+}|N_{\rho,\sigma}^{0+}}=
\deleted{\unsure{M\times M}}\operatorname{Ker}(}
&M\times M\arrow[r]
&M_0)\deleted{\unsure{\longrightarrow0}}\arrow[d,phantom,"\cap" description]\\
&&M
\end{tikzcd}}
\diagram{\begin{tikzcd}[column sep=small]
\deleted{\unsure{N}}\arrow[r,mapsto]
&\deleted{\unsure{\alpha_0,\beta_0}}\;\deleted{\unsure{(\alpha_0,\beta_0)}}
\end{tikzcd}}
\diagram{\begin{tikzcd}[column sep=large]
{(\underset{\substack{\Vert\\\inserted{a_0\lambda_0+a_1\lambda_1}}}{\alpha_0},
\underset{\substack{\Vert\\\inserted{b_0\lambda_0+b_1\lambda_1}}}{\beta_0})}\arrow[r]
& (1-\rho_M)(\alpha_0)-(1-\sigma_M)(\beta_0)
\end{tikzcd}}
\[
\underbrace{\hphantom{(1-\rho_M)(\alpha_0)-(1-\sigma_M)(\beta_0)}}
\]
\inserted{kernel consisting of the $(\alpha_0,\beta_0)$ such that
\[
\begin{aligned}
c_0&\overset{\mathrm{dfn}}{=}a_0+a_1+(b_1-b_0)=0\\
c_1&\overset{\mathrm{dfn}}{=}2a_1-a_0-(b_1-b_0)=0
\end{aligned}
\qquad\bigg|\ \text{cf. }(61)
\]}

Let $\underline N^{0+}$ be the kernel, so that we have
\[
(129)\qquad\underline N_{\rho,\sigma}^{0+}\simeq\underline N^{0+}.
\]

\uline{Let us consider} the homomorphism
\diagram{\begin{tikzcd}
{(130)}\arrow[d,phantom,"\Bigg\{" description]
&\mathbb Z\arrow[r]&N^{0+}\\
{} &n\arrow[r,mapsto]&(n\lambda_0,n\lambda_0)
\end{tikzcd}}

the image \deleted{is} is interpreted as \unsure{made up|made up} of the systems $(a_0,a_1,b_0,b_1)$ such that \deleted{$\unsure{a_0=b_0}$}
\[
a_1=b_1=0,\qquad a_0=b_0,\qquad\unsure{''}
\]
\deleted{whence therefore \unsure{a} \illegible{2}}
\inserted{i.e. we have an exact sequence}
\diagram{\begin{tikzcd}
\deleted{\unsure{(131)}}&
0\arrow[r]&\mathbb Z\arrow[r]&N^{0+}\arrow[r,thick]&\mathbb Z^3\\
&&&(a_0,a_1,b_0,b_1)\arrow[r,mapsto]&(a_1,b_1,a_0-b_0)
\end{tikzcd}}
hence \deleted{\unsure{this}} $N^{0+}/\mathbb Z$ is \deleted{\unsure{flat|free}} a free $\mathbb Z$-module \inserted{of rank $1$},